\chapter{Paquete de plantilla.}

\section{Ejemplos.}



%\textthinmono{Texto en monospace thin.}



\nkLatexPackage{HOla mundo.}¸

\pkg{alias}.

\nkLatexPackage{HOla mundo.}¸

\pkg{alias}.

\begin{nkObjective}
	\begin{itemize}
		\item Objetivo 1.
		\item Objetivo 2.
	\end{itemize}
\end{nkObjective}


%
%\tmplFile{Paquete}.
%
%\tmplCodeItem{Item}.
%
%\tmplBranch{Rama}.
%
%\tmplCodePath{./ruta/ruta1}.
%
%\tmplPath{./ruta/ruta1}.
%
%%\tmplTemporal{Temporal}.
%
%\tmplTerm{Term}.
%
%\tmplTermEn{Term}.
%
%\begin{tmplObjective}	
%	\begin{itemize}
%		\item Describir cómo las grandes compañíias están adoptando DevOps.
%		\item Reconocer que DevOps requiere un cambio en la cultura.
%	\end{itemize}
%\end{tmplObjective}
%
%Para poder adoptar la cultura DevOps, se debe de desaprender la cultura actual.
%
%\begin{tmplSummary}
%	DevOps no se trata de herramientas, se trata de confianza, transparencia, comunicación, coordinación y disciplina en los equipos.
%\end{tmplSummary}
%
%\begin{tmplChapterSummary}
%	\begin{itemize}
%		\item Technology is the enabler of innovation, rather than the driver of innovation. You must have an innovative business idea to leverage technology.
%		\item In 2009, John Allspaw described an innovative approach to managing development and operations that enabled Flickr to complete over ten deploys per day, when many companies were completing fewer than one deploy every six months. This was a key moment in the growth of DevOps.
%	\end{itemize}
%\end{tmplChapterSummary}
%
%\begin{tmplquote}{Patrick Debois, 2009}
%	The term (development and operations) is an extension of agile development environments that aims to enhance the process of software delivery as a whole.
%\end{tmplquote}
%
%\begin{table}[H]
%	\centering
%	\begin{tblr}{
%			width = \textwidth,
%			colspec = {|c|c|X[l,m]|}, % Definimos la alineación en la especificación de columnas
%			hlines,
%			vlines,		
%			rows = {m},
%			row{odd} = {bg=gray!30},
%			%row{even} = {bg=red!30},
%			row{1} = {font=\bfseries, bg=gray!60, c}, % La primera fila centrada
%		}	
%		Label &	Color & Descripción  \\
%		\codeinline{type: feature} & \lblFeature \hspace{5px} \#0E8A16 & Nueva funcionalidad.\\	
%		\codeinline{type: bug} 		& \lblBug \hspace{5px} \#D73A4A		& Algo no funciona correctamente.\\
%		\codeinline{type: hotfix} 	& \lblHotfix  \hspace{5px} \#B60205	& Bug crítico en producción.\\
%		\codeinline{type: docs} 	& \lblDocs \hspace{5px} \#0075CA 		& Documentación.\\
%		\codeinline{type: refactor} & \lblRefactor  \hspace{5px} \#FBCA04	& Mejora de código sin cambio de funcionalidad.\\
%		\codeinline{type: chore} 	& \lblChore \hspace{5px} \#D4C5F9 	& Mantenimiento, dependencias, CI.\\
%	\end{tblr}
%	\caption{\label{tab:workflow-mngmnt_label_worktype} Etiquetas de tipo.}
%\end{table}
%
%

